\begin{lemma}
For every $\eta>0$ there are $\Delta_0$ and $\gamma_0$ such that the
following holds.  Let $G$ be a $\Delta$-regular graph with
$\Delta\ge\Delta_0$, let $\gamma\ge\gamma_0$, and let $I$ be sampled by
\noderef{RandomIndependentSetSampling}.  Fix $r$ and
$X\subseteq N_G(r)$, and assume that no two distinct vertices of $X$ have a
common neighbor outside $N_G[r]$.  If $\mathcal J_2$ and $\mathcal J_3$ are
the independent pairs and independent triples in $X$, then
\[
\mathbb E\binom{|I\cap X|}{2}
-\mathbb E\binom{|I\cap X|}{3}
\ge
\frac{|\mathcal J_2|}{\Delta^2}
-\frac{|\mathcal J_3|}{\Delta^3}
-\eta .
\]
\end{lemma}
\begin{proof}
Write $\overline P_X=\binom{|I\cap X|}{2}$ and
$\overline T_X=\binom{|I\cap X|}{3}$.  For an independent pair
$uv\in\mathcal J_2$, put
$c_{uv}=|N_G(u)\cap N_G(v)|$ and
$\lambda_{uv}=1-c_{uv}/\Delta$.  The clean-neighborhood hypothesis gives
$0<\lambda_{uv}\le1$, because every common neighbor of $u$ and $v$ lies in
$N_G[r]\setminus\{u,v\}$ and so $c_{uv}\le\Delta-1$.

The exact pair identity \noderef{SamplingPairExposureEstimate} gives
\[
\mathbb E\overline P_X
=\sum_{uv\in\mathcal J_2}
\frac{2}{\Delta^2}
\int_{0\le x\le y\le\gamma}
\left(1-\frac{x}{\Delta}\right)^{\lambda_{uv}\Delta}
\left(1-\frac{y}{\Delta}\right)^\Delta\,dy\,dx .
\]
The sampling law in \noderef{RandomIndependentSetSampling} also gives
$0<\gamma$ and, in every non-vacuous case, $\gamma/\Delta\le1$, hence
$\gamma\le\Delta$.  We shall use only the sub-box
$0\le x\le y\le\gamma_0$, so after increasing $\Delta_0$ we may assume this
sub-box lies in $[0,\Delta]^2$.

We need a uniform lower bound for the displayed finite pair integral.  Choose
a fixed number $L$ so large that
\[
2\int_{0\le x\le y\le L} e^{-y}\,dy\,dx
=2\int_0^L y e^{-y}\,dy
> \frac53 .
\]
Then choose a small number $\alpha>0$ so that
$\alpha\le\min\{1,\eta/9\}$,
\[
e^{-\alpha L}\cdot
2\int_0^L y e^{-y}\,dy
> \frac53+4\alpha ,
\qquad
(1-\alpha)^2 e^{-\alpha L}\cdot
2\int_0^L y e^{-y}\,dy
> \frac53+2\alpha .
\]
This is possible because both left sides tend to the first displayed number
as $\alpha\downarrow0$.

The low-overlap regime requires a different fixed window.  After $\alpha$
has been chosen, choose $\gamma_0\ge L$ so large that, uniformly for every
$\lambda\in[\alpha,1]$,
\[
2\int_{0\le x\le y\le\gamma_0} e^{-\lambda x}e^{-y}\,dy\,dx
\ge
\frac{2}{1+\lambda}-\alpha .
\]
Indeed the corresponding integral over the infinite chamber is
\[
2\int_{0\le x\le y<\infty} e^{-\lambda x}e^{-y}\,dy\,dx
=\frac{2}{1+\lambda},
\]
and the tails converge to zero uniformly on the compact parameter interval
$[\alpha,1]$.

Finally choose $\Delta_0$ so large that $\gamma_0\le\Delta$ and the following
finite-product estimates hold for all $\Delta\ge\Delta_0$.  First, uniformly
for $0\le t\le\gamma_0$ and $\lambda\in[\alpha,1]$,
\[
\left(1-\frac t\Delta\right)^{\lambda\Delta}
\ge (1-\alpha)e^{-\lambda t}
\quad\text{and}\quad
\left(1-\frac t\Delta\right)^\Delta\ge(1-\alpha)e^{-t}.
\]
Second, uniformly for $0\le t\le L$ and $0\le\lambda\le\alpha$,
\[
\left(1-\frac t\Delta\right)^{\lambda\Delta}
\ge (1-\alpha)e^{-\lambda t}
\ge (1-\alpha)e^{-\alpha L},
\qquad
\left(1-\frac t\Delta\right)^\Delta\ge(1-\alpha)e^{-t}.
\]
These are standard uniform convergence statements on fixed compact
rectangles, together with $e^{-\lambda t}\ge e^{-\alpha L}$ on
$[0,L]\times[0,\alpha]$.

If $\lambda_{uv}\ge\alpha$, the preceding estimates on
$[0,\gamma_0]$ and the choice of $\gamma_0$ imply that the pair contribution
is at least
\[
\frac1{\Delta^2}(1-\alpha)^2
\left(\frac{2}{1+\lambda_{uv}}-\alpha\right).
\]
Put $A=2/(1+\lambda_{uv})$.  Since $0<\lambda_{uv}\le1$, we have
$1\le A\le2$.  Also $0<\alpha\le1$, and hence
\[
A-(1-\alpha)^2(A-\alpha)
=A(2\alpha-\alpha^2)+\alpha(1-\alpha)^2
\le 2\cdot 2\alpha+\alpha=5\alpha\le6\alpha .
\]
Thus in the case $\lambda_{uv}\ge\alpha$ the pair contribution is at least
\[
\frac1{\Delta^2}
\left(\frac{2}{1+\lambda_{uv}}-6\alpha\right).
\]
If instead $\lambda_{uv}<\alpha$, we use only the
subtriangle $0\le x\le y\le L$.  On this subtriangle the finite product is at
least
\[
(1-\alpha)^2 e^{-\alpha L}e^{-y}.
\]
Therefore the pair contribution is at least
\[
\frac1{\Delta^2}(1-\alpha)^2 e^{-\alpha L}
2\int_0^L y e^{-y}\,dy
\ge \frac1{\Delta^2}\left(\frac53+2\alpha\right),
\]
by the initial choice of $\alpha$.  In both regimes, after this initial
choice of $\alpha$ and the final enlargement of $\Delta_0$, we have the
uniform per-pair bound
\[
\Pr[u,v\in I]\ge
\frac1{\Delta^2}
\left[
1+\frac13\lambda_{uv}
\left(\frac{2}{\lambda_{uv}(1+\lambda_{uv})}-1\right)
\right]
-\frac{6\alpha}{\Delta^2}.
\]
For $\lambda_{uv}\ge\alpha$, this follows from
\[
\frac{2}{1+\lambda}
-\left(1+\frac13\lambda\left(\frac{2}{\lambda(1+\lambda)}-1\right)\right)
=\frac{(1-\lambda)^2}{3(1+\lambda)}\ge0.
\]
For $\lambda_{uv}<\alpha$, the bracketed target is at most $5/3$, while the
fixed-window high-overlap lower bound is at least
$(5/3+2\alpha)/\Delta^2$, which is certainly at least the bracketed target
minus $6\alpha$, divided by $\Delta^2$.
Finally, $|\mathcal J_2|\le\binom{|X|}{2}\le\Delta^2/2$, because
$X\subseteq N_G(r)$ and $G$ is $\Delta$-regular.  Therefore the summed loss
from the terms $6\alpha/\Delta^2$ is at most $3\alpha\le\eta/3$.

Summing over all pairs gives
\[
\mathbb E\overline P_X\ge
\frac{|\mathcal J_2|}{\Delta^2}
+\frac1{3\Delta^2}\sum_{uv\in\mathcal J_2}
\lambda_{uv}
\left(\frac{2}{\lambda_{uv}(1+\lambda_{uv})}-1\right)
-\eta/3 .
\]

Now apply \noderef{SamplingTripleExposureEstimate} with error budget
$\eta/3$.  It gives
\[
\mathbb E\overline T_X\le
\frac{|\mathcal J_3|}{\Delta^3}
+\frac1{3\Delta^3}
\sum_{uvw\in\mathcal J_3}\sum_{ab\in\{uv,uw,vw\}}
\left(\frac{2}{\lambda_{ab}(1+\lambda_{ab})}-1\right)
+\eta/3 .
\]
For a fixed independent pair $uv$, any vertex $w$ for which $uvw$ is an
independent triple lies in $X\subseteq N_G(r)$ and is not a common neighbor
of $u$ and $v$, because it is adjacent to neither $u$ nor $v$.  Therefore
there are at most
$\Delta-c_{uv}=\lambda_{uv}\Delta$ possible third vertices for this fixed
pair.  Re-indexing the double sum by its distinguished pair gives
\[
\sum_{uvw\in\mathcal J_3}\sum_{ab\in\{uv,uw,vw\}}
\left(\frac{2}{\lambda_{ab}(1+\lambda_{ab})}-1\right)
\le
\sum_{uv\in\mathcal J_2}\lambda_{uv}\Delta
\left(\frac{2}{\lambda_{uv}(1+\lambda_{uv})}-1\right).
\]
Consequently
\[
\mathbb E\overline T_X\le
\frac{|\mathcal J_3|}{\Delta^3}
+\frac1{3\Delta^2}\sum_{uv\in\mathcal J_2}
\lambda_{uv}
\left(\frac{2}{\lambda_{uv}(1+\lambda_{uv})}-1\right)
+\eta/3 .
\]

Subtracting this upper bound from the preceding lower bound for
$\mathbb E\overline P_X$ cancels the whole overlap-dependent correction and
leaves
\[
\mathbb E\overline P_X-\mathbb E\overline T_X
\ge
\frac{|\mathcal J_2|}{\Delta^2}
-\frac{|\mathcal J_3|}{\Delta^3}
-2\eta/3 .
\]
Since $\eta>0$, this implies the stated inequality with final error
$\eta$.
\end{proof}
